\documentclass{article}
\usepackage[T1]{fontenc}
\usepackage{lmodern}
\usepackage{graphicx}
\usepackage{amsmath,amssymb}
\usepackage[a4paper,margin=2cm]{geometry}
\usepackage[absolute,overlay]{textpos}
\setlength{\TPHorizModule}{1mm}
\setlength{\TPVertModule}{1mm}
\usepackage[indonesian]{babel}
\usepackage{hyperref}
\setlength{\emergencystretch}{2em}
% unit: Halaman 2

\begin{document}

% === Halaman 2 (python/native) ===

Latar Belakang

• DES dianggap sudah tidak aman karena kunci dapat ditemukan secara 

brute-force.

• Perlu diusulkan standard algoritma baru sebagai pengganti DES.

• National Institute of Standards and Technology  (NIST) mengusulkan 

kepada Pemerintah Federal AS untuk sebuah standard kriptografi baru.

• NIST mengadakan kompetisi terbuka untuk membuat standard algoritma 

kriptografi yang baru. Standard tersebut kelak diberi nama Advanced 

Encryption Standard (AES).

\end{document}
